\section{Time value, required returns and present value}

\begin{qfdefinition}[Required-return components: real rate, expected inflation and risk premia; the Fisher relation, exact for realised inflation, approximate in its additive form]\label{def:required-return-components}
\qffield{Statement} A \emph{nominal} rate is quoted in money; a \emph{real} rate measures the growth of purchasing power. The Fisher relation, $1+i=(1+R)(1+h)$, links them, with $i$ the nominal rate, $R$ the real rate and $h$ expected inflation: a 6 percent real rate with 2 percent expected inflation requires a nominal rate of 8.12 percent \citep[sec.~7.4]{principles-of-finance-2026}. The same text also uses the subtraction ``real interest rate $=$ nominal interest rate $-$ inflation rate'' (sec.~3.4), which is the first-order approximation of the exact relation. The short-term government borrowing rate is ``often referred to as the risk-free rate of interest'', and lenders ``charge a premium to borrowers who are higher risk'' (sec.~3.4).

A nominal government yield is not just the real rate plus expected inflation. In the Federal Reserve staff's term-structure model, nominal yield $=$ real yield $+$ expected inflation $+$ inflation risk premium, and real yield $=$ expected average future real short rate $+$ real term premium; the two premia are ``extra compensations bond investors demand for bearing inflation risks and real interest rate risks, respectively'' \citep[eqs.~(1)--(2)]{inflation-compensation-decomposition-2019}. Inflation-indexed bonds trade at a yield above the true real yield by a liquidity premium, so their breakeven (inflation compensation) equals expected inflation plus the inflation risk premium minus that liquidity premium (eqs.~(3)--(4)). A borrower that may default pays a credit spread over the risk-free rate \citep[chap.~4, sec.~4.1, p.~98]{s7-hulloptions}.

\qffield{Notation} $R$ real rate, $h$ inflation, $i$ nominal rate, all for the same horizon and compounding. Interpretation: with \emph{realised} inflation $h$, or within one specified inflation scenario, $R=(1+i)/(1+h)-1$ is an identity, the realised (ex post) real return, and $i=R+h+Rh$ exactly. Approximation: $i\approx R+h$, with error $Rh$; with continuously compounded rates, $\ln(1+i)=\ln(1+R)+\ln(1+h)$ holds exactly. Interpretation, a decomposition by convention: the required yield on a bond or the required return on an equity claim $=$ nominal risk-free yield of matching maturity $+$ premia for the claim's risks (credit or default, liquidity, and for equity an equity risk premium). Premia stated as spreads add only as an approximation to products of growth factors.

\qffield{Conditions} When a nominal rate is agreed, future inflation is unknown, so there are two real rates: the \emph{ex ante} real rate expected when the loan is made, $i-E\pi$, and the \emph{ex post} real rate actually realised, $i-\pi$; the nominal rate ``can adjust only to expected inflation'', so the Fisher effect is written $i=r+E\pi$ \citep[chap.~5, sec.~5-3]{macroeconomics-2019}. These are the additive forms. \emph{Uncertain inflation} (Interpretation, derived here): deflating a nominal rate by expected inflation is not the expected real return, because the realised real return $(1+i)/(1+h)-1$ is a convex function of $h$, so its mean exceeds its value at the mean (Jensen's inequality). With a fixed nominal return of 10 percent and inflation of 0 or 20 percent with equal probability, deflating by mean inflation, 10 percent, gives 0 percent, while the mean realised real return is $\tfrac12\bigl(10+(1.10/1.20-1)\times100\bigr)=0.8333$ percent. A required nominal yield under uncertainty therefore also contains compensation for inflation risk, which is what the stochastic decomposition below separates out; the teaching calculation in Practice assumes a single known inflation rate. The premia in the Federal Reserve decomposition are model estimates from an affine no-arbitrage term-structure model of U.S.\ Treasury and TIPS yields, not observed prices; the 2019 note updates estimates published for 1990--2013 (Kim, Walsh and Wei, introduction and model sections). Interpretation: credit, liquidity and equity premia are compensation the market requires; a component cannot be read from one yield without a model or a matched comparison, such as a breakeven or a credit spread over the same-maturity government yield.

\qffield{Practice} Example, original, with inflation treated as known: a real rate of 1.5 percent and inflation of 2.5 percent give a nominal rate of $(1.015)(1.025)-1=4.0375$ percent; the additive approximation, 4.0 percent, misses $Rh=0.0375$ percentage points. Adding an inflation risk premium of 0.30 points gives a nominal risk-free yield of 4.3375 percent. A corporate bond with a credit spread of 1.20 and a liquidity premium of 0.30 points then needs 5.8375 percent, and an equity claim with an equity risk premium of 5.0 points needs 9.3375 percent. The premia here are illustrative, not estimates.

\qffield{Pitfalls} The breakeven inflation rate is not the expected inflation rate; it also contains the inflation risk premium and, with the opposite sign, the indexed bond's liquidity premium (eq.~(4)). Subtracting inflation from a nominal rate is an approximation; at high inflation the error is material: at 20 percent inflation and a 10 percent nominal rate, the exact real rate is $-8.33$ percent, the approximation $-10$ percent. Do not mix compounding conventions when adding components.

\qffield{See also} \hyperref[def:return-measures]{return measures} (real return); \hyperref[def:tips-breakevens]{inflation-linked bonds and breakevens}; \hyperref[mod:forward-rates-fra]{forward rates} (term-structure theories); \hyperref[meth:instrument-present-value]{present value of instrument cash flows}.

\end{qfdefinition}

\begin{qfmethod}[Present value of bond and equity cash flows, and the implied yield, required return and growth rate that equate price to present value]\label{meth:instrument-present-value}
\qffield{Statement} A bond's price is the present value of its coupons plus the present value of its face value, discounted at the yield to maturity per period; for semiannual coupons the annual coupon and the yield, quoted as a nominal annual rate with semiannual compounding, are halved and the number of periods doubled. The yield cannot be isolated algebraically and is found with a financial calculator or spreadsheet; a bond priced above par has a coupon above its yield, one below par a coupon below its yield \citep[sec.~10.2]{principles-of-finance-2026}. A bond can equally be priced by discounting each cash flow at the zero (spot) rate for its own maturity; its yield is the single rate that reproduces that price \citep[chap.~4, secs.~4.5--4.6, pp.~104--105]{s7-hulloptions}.

A share is valued by the present value of its expected dividends. With a constant dividend $D$ forever and required return $r$, the value is the perpetuity $D/r$; with dividends growing at a constant rate $g$, it is $D_0(1+g)/(r-g)=D_1/(r-g)$, which ``requires the return rate $r$ to be greater than the growth rate $g$'' \citep[sec.~11.2]{principles-of-finance-2026}. The same equation, given a price, can be solved for the growth it implies (sec.~11.2, Example~2).

\qffield{Notation} $C$ coupon per period, $F$ face value, $y$ yield per period, $N$ periods, $z_t$ zero rate for maturity $t$, annual compounding and cash flows at period ends:
\[ P=\sum_{t=1}^{N}\frac{C}{(1+y)^t}+\frac{F}{(1+y)^N}=C\,\frac{1-(1+y)^{-N}}{y}+\frac{F}{(1+y)^N},\qquad P=\sum_{t=1}^{N}\frac{CF_t}{(1+z_t)^t}. \]
Quotation (Interpretation): halving applies to a yield quoted as a nominal annual rate compounded semiannually, $y_{\mathrm{nom}}$, whose half-year rate is $y_{\mathrm{nom}}/2$. A yield quoted as an \emph{effective} annual rate $y_{\mathrm{eff}}$ has half-year rate $(1+y_{\mathrm{eff}})^{1/2}-1$; the two describe the same bond when $(1+y_{\mathrm{nom}}/2)^2=1+y_{\mathrm{eff}}$ (\hyperref[def:interest-rate-measures]{compounding conversions}). Equity: $V_0=D/r$ (constant dividend, as for non-callable fixed-rate preferred stock); $V_0=D_1/(r-g)$, $r>g$; two-stage growth: $V_0=\sum_{t=1}^{n}D_t(1+r)^{-t}+(1+r)^{-n}D_{n+1}/(r-g_L)$, the last term the present value of the terminal value at $n$. Interpretation, derived by rearranging: implied required return $r=D_1/P_0+g$, dividend yield plus growth; implied growth $g=(rP_0-D_0)/(P_0+D_0)$.

\qffield{Conditions} The bond formula values the cash flows on a coupon date; pricing between coupon dates needs accrued interest and day counts (\hyperref[conv:us-rate-quotation]{quotation conventions}). A single yield is a summary of the zero curve, not a rate earned on each cash flow, and realising it requires reinvestment at that yield. The dividend models are valuation identities given their inputs: the expected dividends, the constant $g<r$ forever, and the required return, which itself contains the premia of \hyperref[def:required-return-components]{required-return components}. Value is very sensitive to $r-g$. Implied returns and growth are what the price requires under the model, not forecasts.

\qffield{Practice} Examples, original, annual compounding. A zero-coupon bond paying 1{,}000 in three years at 4 percent is worth 888.9964. A three-year 5 percent annual-coupon bond, face 100, is worth 102.7751 at a 4 percent yield (premium bond, coupon above yield); at zero rates of 3.0, 3.8 and 4.1 percent it is worth 102.5709, and its yield, 4.0723 percent, lies between the zero rates and closest to the three-year rate. Priced at 98.00 the same bond yields 5.7447 percent, found by bisection.

A preferred share paying a fixed 6 a year at a required 7.5 percent is worth 80.00. A share with $D_0=2.00$, $g=4$ percent and $r=9$ percent is worth $2.08/0.05=41.60$, and at that price $2.08/41.60+0.04=9$ percent. With 10 percent growth for two years and 4 percent after, the dividends are 2.20 and 2.42, the terminal value at year 2 is $2.42\times1.04/0.05=50.336$, and the value is 46.4220. At a price of 52 and $r=9$ percent, the implied constant growth is $(0.09\times52-2)/54=4.9630$ percent. Each figure was also obtained by summing the cash flows directly.

\qffield{Pitfalls} The constant-growth formula fails when $g\ge r$. Use the dividend expected next period, $D_1$, not the last one paid. Match the rate to the cash-flow period and compounding: cash flows of 5 in six months and 105 in a year are worth 100.2219 at a 10 percent \emph{effective} annual yield (half-year rate 4.8809 percent), the same as at a nominal 9.7618 percent compounded semiannually; halving the effective 10 percent instead gives 100.00, a mispricing. A higher yield means a lower price for given cash flows.

\qffield{See also} \hyperref[def:interest-rate-measures]{interest-rate measures} (compounding, annuities, zero and par yields); \hyperref[res:cash-flow-additivity]{cash-flow additivity}; \hyperref[def:return-measures]{return measures} (the IRR).

\end{qfmethod}

\begin{qfresult}[Cash-flow additivity and no arbitrage: the value of a sum of cash flows is the sum of their values, applied to implied forward rates, forward exchange rates and a replicated option]\label{res:cash-flow-additivity}
\qffield{Statement} Two portfolios with identical payoffs in every state must have identical values today; ``if this were not the case, an arbitrageur could buy the less expensive portfolio and sell the more expensive one'' \citep[chap.~11, sec.~11.4, p.~255]{s7-hulloptions}. Applied to cash flows dated in the future, value is additive: a stream is worth the sum of its dated cash flows, each discounted at the rate for its date, and a position assembled from simpler instruments is worth the sum of their values. Three consequences follow.

\begin{itemize}
\item \emph{Forward rates.} Investing for $T_2$ years, or for $T_1$ years and then at the forward rate from $T_1$ to $T_2$, must give the same amount; a borrower and lender at the zero rates can lock in the forward rate today: borrowing for one year and investing for two creates an outflow at year 1 and an inflow at year 2, a loan made between those dates at the forward rate, and borrowing for the longer and investing for the shorter term locks in forward borrowing \citep[chap.~4, sec.~4.8, p.~110]{s7-hulloptions}. Rolling short borrowings over at whatever rates prevail later is not a lock-in; it is a bet that rates will differ from the forwards (Business Snapshot~4.1, p.~111).
\item \emph{Forward exchange rates.} Borrowing domestic currency, converting at spot, investing the foreign currency and selling the proceeds forward must cost nothing at maturity, which gives covered interest parity, $F_0=S_0e^{(r-r_f)T}$ with continuous compounding (chap.~5, sec.~5.10, pp.~137--140; \hyperref[res:cost-of-carry-forward-price]{cost of carry}).
\item \emph{Options.} When the underlying can move only to two prices, a position of $\Delta$ shares and a riskless bond can reproduce an option's payoff in both states, so the option is worth that position (chap.~13, sec.~13.1, pp.~288--290); and a call plus a zero-coupon bond paying $K$ has the same payoff as a put plus the share, so $c+Ke^{-rT}=p+S_0$ (sec.~11.4, eq.~(11.6), p.~256).
\end{itemize}

\qffield{Notation} Interpretation, the discrete-compounding forms used in the examples. Zero rates $z_1,z_2$ (annual): $(1+z_2)^2=(1+z_1)(1+f_{1,1})$, so $f_{1,1}=(1+z_2)^2/(1+z_1)-1$. Spot $S_0$ and forward $F_0$ in domestic currency per unit of foreign currency, domestic and foreign rates $i_d,i_f$ for the same period: $F_0=S_0(1+i_d)/(1+i_f)$. One period, share $S_0$ moving to $S_u$ or $S_d$, option payoffs $f_u,f_d$, rate $r$ for the period: $\Delta=(f_u-f_d)/(S_u-S_d)$, bond holding $B=(f_d-\Delta S_d)/(1+r)$ in present value, option value $\Delta S_0+B$; parity $c+K/(1+r)=p+S_0$.

\qffield{Conditions} No arbitrage among participants who can trade the component instruments without transaction costs, borrow and lend at the stated rates and, for parity, hold European options with the same strike and expiry on a share that pays no dividend before expiry (Hull, secs.~5.3 and~11.4). Additivity values claims whose cash flows are fixed or replicable; it does not say what an unreplicable payoff is worth. The two-state option model is a teaching device: real prices take many values, and the replication must be rebalanced (Hull, chap.~13). Forward rates and forwards are prices that can be locked in, not forecasts.

\qffield{Practice} Examples, original. \emph{Forward rate.} With annual zero rates of 3.0 and 3.8 percent, the one-year rate one year forward is $1.038^2/1.03-1=4.6062$ percent. Buying a two-year zero that pays 1{,}000 for 928.1225 and borrowing that amount for one year at 3.0 percent gives dated cash flows of 0 now, $-955.9662$ at year 1 (repaying the loan) and $+1{,}000$ at year 2: a loan of 955.9662 from year 1 to year 2 at $1{,}000/955.9662-1=4.6062$ percent, fixed today. The reverse trades, selling the two-year zero short and lending the 928.1225 for one year, give $+955.9662$ at year 1 and $-1{,}000$ at year 2, borrowing locked at the same rate. Repaying the year-1 debt by a new one-year loan is a different position: its rate is unknown today, and at 8 percent it would cost 1{,}032.4434 at year 2.

\emph{Forward exchange rate.} At 1.1000 USD per EUR, a one-year USD rate of 4.0 percent and EUR rate of 2.5 percent give $1.1000\times1.04/1.025=1.116098$ USD per EUR, a forward premium of 160.98 pips. Borrowing USD 1{,}100, buying EUR 1{,}000, depositing to EUR 1{,}025 and selling that forward returns USD 1{,}144, the amount owed.

\emph{Option.} A share at 50 moves to 60 or 40, and the period rate is 5 percent. A call struck at 50 pays 10 or 0, so $\Delta=0.5$ and the bond holding is $-20/1.05=-19.0476$: buy half a share and borrow 19.0476, at a value of 5.9524. The position pays 10 or 0, as the call does. Parity gives the put struck at 50 as $5.9524+50/1.05-50=3.5714$, which a separate replication of the put confirms.

\qffield{Pitfalls} Additivity requires the same date and the same currency: amounts must be discounted to a common date before they are added. A forward exchange rate quoted the other way round (foreign per domestic) inverts the formula. The replicating $\Delta$ and bond change when the share price or time changes; a static position replicates only the one-period payoff.

\qffield{See also} \hyperref[mod:forward-rates-fra]{forward rates and FRAs}; \hyperref[res:cost-of-carry-forward-price]{cost-of-carry forward price}; \hyperref[res:synthetic-short-put-call]{synthetic short through options} (parity with dividends); \hyperref[meth:instrument-present-value]{present value of instrument cash flows}.

\end{qfresult}
